\section{The K-theory transfer: source distinctions and new exact tests}
\label{sec:ktheory}

\subsection{Why this is beyond a usual RH consequence}

The groups $K_{4m}(\Z)$ encode higher algebraic information about the
integers. Their proposed vanishing is connected to cyclotomic class groups
and strict Galois-cohomology groups. A bound on one of those groups inside
the original RH proof can therefore have an application that does not
factor through the final assertion about zeta zeros.

Mok--Zheng's November 24 preprint claims
\[
 K_{4m}(\Z)=0\qquad(1\le m\le30).
\]
Pintar's current preprint states a conditional range through degree $104$,
retaining the original Step 17 arithmetic construction as a hypothesis.
The former reports that its modified Gaussian Lean entry point has not
been fully rebuilt; the latter invokes a common-representation extraction
from a separate research archive. These are different proof routes and
different verification statuses \cite{mokzheng,pintar}.

The earlier ProveIt analysis already recognized these developments.
The useful new work is to expose the interfaces and test arithmetic
beyond the published ranges, while keeping the unverified logical
steps visible.

\subsection{The short repository does not export the required strict bound}

The current short source proves
\path{New.exists_length_shaOdd_le}. For a fixed cyclotomic CM field
$E\subseteq\Q(\mu_N)$ and a fixed range of Tate exponents, it supplies
one nonzero integer whose valuations bound lengths of certain odd
cocycle modules. Its documentation explicitly distinguishes that
model from ordinary $\Sha^1$ \cite{oddsha}.

The newest predecessor supplies a comparison with an ordinary local
cohomology kernel, for odd coefficient primes and positive Tate
exponents. Its conjugation eigenspace is
\[
 \Sha^1(E,(F/\mathcal O)(-e))^{c=(-1)^{e+1}}.
\]
For even $e$, this is the natural minus eigenspace. Rational classes
needed in the direct $K_{4m}(\Z)$ route lie in the natural plus
eigenspace after restriction. A length bound on the former is not a
vanishing theorem for the latter. Also, an integer obtained after
fixing $E$ cannot be used uniformly after replacing $E$ by
$\Q(\mu_\ell)$ as $\ell$ varies \cite{priorbeyond}.

The original interface used by the K-theory papers is instead
\path{UpperTheory.V.Z.J.OneCharU_s}, in the much larger original
development, with the caller \path{SmallClassNumbers/CorGKOf.lean}.
Those names are absent from the shorter source tree. We inspected
the original definitions at the separately pinned mirror revision
\code{7cbf7a0c32fca1badf05931879a37c3f8a2b03b4}
\cite{originalinterface}.

In the pointwise formulation, a Galois element is selected first;
the finite family of admissible automorphic witnesses may depend
on that element. Pintar's stronger route extracts a common
representation from the construction behind the inequality.
The inequality alone does not exchange these quantifiers.

There is a further concrete field issue. For the original parameters
$N=2$, an even exponent, and trivial character, the fields
$\Q(\sqrt{-3})$ and $\Q(\sqrt{-7})$ satisfy the relevant
odd-discriminant coprimality condition. The Gaussian field $\Q(i)$
does not: $\gcd(4,2)\ne1$. Using it therefore requires the additional
proof transfer described by Mok--Zheng, rather than substitution
into the unmodified original hypothesis.

\subsection{How to aggregate finite certificates correctly}

For a finite set of valid Frobenius probes $q$, suppose a nonzero
Selmer group forces at least one of the integral branch determinants
$A_{qj}$ to be divisible by a coefficient prime $\ell$.
Then
\[
 \bigl(\forall q\ \exists j:\ell\mid A_{qj}\bigr)
 \quad\Longleftrightarrow\quad
 \ell\mid\gcd_q\prod_j A_{qj}.
\]
This is the product-before-gcd rule already proved in the dated
predecessor. Validity of the probe, integrality of the chosen lattice,
and the target trace normalization are hypotheses of the
divisibility statement, not consequences of the integer calculation.

There is an additional useful refinement for \emph{stronger} inputs.

\begin{proposition}[Exact certificate for alternative length bounds]
\label{prop:lcm}
Let the probe set and each corresponding branch set be finite and
nonempty. Let $A_{qj}$ be nonzero integers, and suppose a nonnegative integer
$r_\ell$ satisfies
\[
 \forall q\ \exists j:\quad r_\ell\le v_\ell(A_{qj}).
\]
Set
\[
 C=\gcd_q\operatorname{lcm}_j|A_{qj}|.
\]
Then
\[
 r_\ell\le v_\ell(C)
       =\min_q\max_jv_\ell(A_{qj}).
\]
The right side is the largest bound permitted by these numerical
inequalities alone.
\end{proposition}

\begin{proof}
At a fixed probe the existence of a branch is equivalent to
$r_\ell\le\max_jv_\ell(A_{qj})$. Intersecting the bounds over probes
takes the minimum. Prime valuations turn lcm into maximum and gcd
into minimum. Taking $r_\ell$ equal to that minimum demonstrates
sharpness as a statement about the numerical hypotheses.
\end{proof}

If $r_\ell$ is the length of a finite $\Z_\ell$-module, its
$\ell$-primary exponent is consequently at most $\ell^{v_\ell(C)}$.
This conclusion requires the displayed length hypothesis. The
mere nonzero-witness condition imposes no such bound: for rows
$(2,3)$ and $(3,2)$, the prime $2$ survives all probes regardless
of whether the module is $\Z/2\Z$ or $\Z/2^{100}\Z$.

If the genuine source estimate is a \emph{sum} of valuations of
simultaneous witnesses, products are still needed. The lcm
refinement must not be silently applied to Mok--Zheng's valuation
sum. It describes a precise improvement available only after a
stronger branchwise length theorem has been established.

\subsection{A finite arithmetic theorem beyond the cited degree range}

For $d\in\{3,7\}$ let $\chi_{-d}(a)=(a/d)$, an odd quadratic
Dirichlet character. Define generalized Bernoulli numbers by
\[
 \sum_{m\ge0}B_{m,\chi_{-d}}\frac{t^m}{m!}
 =\frac{t\sum_{a=1}^{d-1}\chi_{-d}(a)e^{at}}{e^{dt}-1},
\]
and put
\[
 b_d(n)=\frac{B_{n+1,\chi_{-d}}}{n+1},\qquad
 A_d(e)=\left|\num b_d(e)\,\num b_d(2e)\right|.
\]
Numerators and denominators are always reduced to lowest terms.

\begin{theorem}[Exact two-field unit screening]\label{thm:screening}
For every even $e$ with $54\le e\le80$ and every prime $\ell>109$,
there is $d\in\{3,7\}$ for which
\[
 b_d(0),\quad b_d(e),\quad b_d(2e)
\]
are all $\ell$-adic units. In the two cases $e=62,64$, the same
conclusion holds for every prime $\ell\notin\{2,3,7\}$.
\end{theorem}

The key data for the next two degrees are:

\begin{center}
\begin{tabular}{@{}rrrrr@{}}
\toprule
$e$ & $2e$ & Denominators for $d=3$ &
Denominators for $d=7$ & $\gcd(A_3,A_7)$\\
\midrule
62 & 124 & $(27,3)$ & $(49,1)$ & $4$\\
64 & 128 & $(3,9)$ & $(1,7)$ & $4$\\
\bottomrule
\end{tabular}
\end{center}

Here each pair gives the denominators of $b_d(e),b_d(2e)$.
The full screening table is:

\begin{center}
\begin{tabular}{@{}rrl@{\qquad}rrl@{}}
\toprule
$e$ & $2e$ & $\gcd(A_3,A_7)$ & $e$ & $2e$ & $\gcd(A_3,A_7)$\\
\midrule
54&108&$4\cdot37\cdot109$&68&136&$4$\\
56&112&$4$&70&140&$4$\\
58&116&$4\cdot47$&72&144&$4\cdot37^2$\\
60&120&$4$&74&148&$4$\\
62&124&$4$&76&152&$4$\\
64&128&$4$&78&156&$4\cdot79^2$\\
66&132&$4\cdot67^2$&80&160&$4$\\
\bottomrule
\end{tabular}
\end{center}

\begin{proof}
The supplied exact program makes 84 rational-value comparisons
in two ways. First, comparison of coefficients after multiplication
by $e^{dt}-1$ gives the recurrence
\[
 b_d(n)=
 \frac{\displaystyle
 \sum_{a=1}^{d-1}\chi_{-d}(a)a^{n+1}
 -\sum_{j=2}^{n+1}\binom{n+1}{j}d^j b_d(n+1-j)}
 {(n+1)d}.
\]
Second, ordinary Bernoulli polynomials give the independent formula
\[
 b_d(n)=\frac{d^n}{n+1}
 \sum_{a=1}^{d-1}\chi_{-d}(a)
 \sum_{j=0}^{n+1}\binom{n+1}{j}B_j(a/d)^{n+1-j}.
\]
Both methods use integers and rational fractions and agree exactly.
All computed denominators have prime factors only in $\{3,7\}$.
For every table entry the output supplies a B\'ezout equality
$uA_3+vA_7=g$ with the displayed $g$, together with exact divisibility
of $A_3,A_7$ by $g$. These certify each gcd without factoring the
large numerator products.

For a prime $\ell>109$, the table implies that $\ell$ cannot divide
both numerator products. Select the field whose product is prime
to $\ell$. The denominators are also prime to $\ell$.
Finally $b_3(0)=-1/3$ and $b_7(0)=-1$. For $e=62,64$, the common
gcd is $4$, so the same argument only excludes $2,3,7$.
\end{proof}

This is a finite arithmetic theorem, with a short exact certificate
checker. Its role in the larger program needs care. The three unit
conditions belong to the common-representation and Bernoulli-exclusion
route described by Pintar. They do not, by themselves, remove all
branches from the weaker pointwise-witness route. They also give no
unit bound on the required Hecke determinant.

The screening range is compatible with the established cyclotomic
computation of Hart--Harvey--Ong, which verifies Vandiver for primes
below $2^{31}$ \cite{hho}. That external result handles small
coefficient primes only after the relevant K-theoretic comparison
has been justified. No new Vandiver verification is claimed here.

\subsection{Exact dimensions of the next comparison spaces}

These calculations also quantify the cost of the next arithmetic
step. Let $P_n$ be the space over $\C$ of polynomials homogeneous
of degree $n$ in each of the two triples $x,y$. Put
\[
 q=x_1y_1+x_2y_2+x_3y_3,\qquad W_n=P_n/qP_{n-1}.
\]
The group $\Gamma_h=\mu_h^3\rtimes S_3$ scales $x_i$ and $y_i$
inversely and permutes the coordinates. The swap
$J:P(x,y)\mapsto P(y,x)$ normalizes this action.

Let $\mathcal O_h(n)$ be the set of unordered triples of pairs
$(a_i,b_i)$ with $\sum a_i=\sum b_i=n$ and
$a_i\equiv b_i\pmod h$. Let $C_h(n)$ count those triples fixed
after swapping the entries of each pair.

\begin{proposition}[Orbit formula for the signed dimension]
\[
 \begin{split}
 d_h(n)&=|\mathcal O_h(n)|-|\mathcal O_h(n-1)|,\\
 t_h(n)&=C_h(n)-C_h(n-1),\\
 \dim(W_n^{\Gamma_h})^-&=\frac{d_h(n)-t_h(n)}2.
 \end{split}
\]
\end{proposition}

\begin{proof}
The coordinate congruences specify the invariant monomials under
$\mu_h^3$. Orbit sums give a basis of their $S_3$ invariants.
The swap permutes that basis, so its trace is the number of fixed
orbits. Multiplication by $q$ is injective and equivariant.
Taking invariants is exact in characteristic zero by averaging
over the finite group. Subtraction gives the dimension and trace
on the quotient. The final identity follows from $J^2=1$.
\end{proof}

With $n=e-1$, exact enumeration gives:

\begin{center}
\begin{tabular}{@{}rrrrrr@{}}
\toprule
$2e$ & $h$ & $n$ & Full dimension & Swap trace & Negative dimension\\
\midrule
124&6&61&1110&10&550\\
128&6&63&1221&11&605\\
124&4&61&2490&10&1240\\
128&4&63&2731&11&1360\\
\bottomrule
\end{tabular}
\end{center}

The stabilizers $h=6$ and $h=4$ correspond to standard
Eisenstein and Gaussian ternary lattices. The formula computes
these finite representation spaces unconditionally. Identifying
them with an exhaustive automorphic comparison space still uses
the relevant field's global and local transfer.

For comparison, the odd-discriminant $\Q(\sqrt{-7})$ route uses
two lattice classes. Its standard-class negative ranks in the
signed elimination formula are $\binom{32}{3}=4960$ and
$\binom{33}{3}=5456$. Thus the Gaussian route offers much smaller
second-field calculations, but only after its modified field
hypothesis has been validated. This is a concrete computation
versus proof-transfer tradeoff, not evidence for a shortcut.

\subsection{Independent reconstruction of two small Hecke operators}

To test more than stored integer outputs, we rebuilt two low-weight
operators from the published 18-neighbor formula. For each coordinate
pair $i<j$ and sixth root of unity $u$, use the identity on the
remaining coordinate and
\[
 g_{i,j,u}=\frac12
 \begin{pmatrix}1&b\\-\bar b&1\end{pmatrix},
 \qquad b=\sqrt{-3}\,u.
\]
The checker verifies unitarity and pairwise coset inequivalence
at $2$. It evaluates the invariant basis on six fresh incidence
pairs satisfying $x\cdot y=0$ and solves exact nonsingular
interpolation systems. No stored operator matrix is used as input.

To specify the matrices, put $z_i=x_i y_i$,
$s_2=\sum_{i<j}z_i z_j$ and $s_3=z_1z_2z_3$.
The ordered basis for $e=6$ is $(s_2s_3)$, and for $e=8$ it is
\[
 \left(s_2^2s_3,\quad
       \sum_{i\ne j}(x_i y_j)^6z_i,\quad
       \sum_{i\ne j}(x_i y_j)^6z_j\right).
\]
Columns give the images of these ordered basis elements.
The reconstructed values are
\[
 T_2(e=6)=\frac{99}{128},\qquad
 2048\,T_2(e=8)=
 \begin{pmatrix}
 10215&-81081&-81081\\
 -81&-2421&11907\\
 -81&11907&-2421
 \end{pmatrix}.
\]
For
\[
 D_e=\det\left(4^{e-1}(T_2+3I)
                   -(4^{2e}+4^e+1)I\right)
\]
the exact outputs are
\[
 D_6=-16777449=-3^5\cdot13\cdot47\cdot113,
 \qquad
 D_8=-79228276291965330830049420825.
\]
The complete factorization of $|D_8|$ is
\[
 3^{13}5^2\,7\,11^2\,23\,401\,809\,1753\,13109\,13687.
\]
All displayed factors and multiplications are checked exactly.
The completeness and interpretation of the geometric neighbor
model remain the published mathematical input \cite{pintarcode}.

Separately, the upstream arithmetic verifier was rerun on its
stored large certificate bundles, and the original small two-probe
generators were rerun. Those checks passed. Replaying a stored
large determinant is weaker evidence than regenerating the
large matrix; our new standalone reconstruction only addresses
the two small compact operators just displayed.
